% ============================================================
%  TEMPLATE PROGETTO - Software Security e Blockchain
%  UNIVPM - Laurea Magistrale in Ing. Informatica e dell'Automazione
%
%  Compilazione consigliata:  latexmk -pdf main.tex
%  (oppure carica la cartella su Overleaf: main.tex è il file principale)
% ============================================================
\documentclass[11pt,a4paper,oneside]{report}

% ---------- Lingua e codifica ----------
\usepackage[utf8]{inputenc}
\usepackage[T1]{fontenc}
\usepackage[italian]{babel}
\usepackage{lmodern}
\usepackage{microtype}

% ---------- Layout ----------
\usepackage[a4paper, top=3cm, bottom=3cm, left=3cm, right=2.5cm]{geometry}
\usepackage{setspace}
\onehalfspacing
\usepackage{parskip}
\setlength{\parskip}{0.6em}
\usepackage{fancyhdr}
\usepackage{titlesec}
\usepackage{enumitem}

% ---------- Colori e grafica ----------
\usepackage[dvipsnames,svgnames]{xcolor}
\definecolor{navy}{HTML}{0B1F3A}
\definecolor{teal}{HTML}{14A3A3}
\definecolor{codebg}{HTML}{F4F7FA}
\definecolor{codekw}{HTML}{0B5FA5}
\definecolor{codecom}{HTML}{5A7D5A}
\definecolor{codestr}{HTML}{B3541E}
\usepackage{graphicx}
\usepackage{tikz}
\usetikzlibrary{arrows.meta, positioning, shapes.geometric, calc, fit, backgrounds}

% ---------- Tabelle, figure, matematica ----------
\usepackage{booktabs}
\usepackage{tabularx}
\usepackage{array}
\usepackage{caption}
\usepackage{subcaption}
\usepackage{amsmath, amssymb}
\captionsetup{font=small, labelfont={bf,color=navy}}

% ---------- Codice (Solidity, JSON, Python, shell) ----------
\usepackage{listings}
\lstdefinelanguage{Solidity}{
  keywords={pragma, solidity, contract, function, public, private, internal, external,
            view, pure, returns, return, mapping, address, uint256, uint, bool, string,
            bytes32, bytes, event, emit, modifier, require, msg, constructor, memory, storage,
            struct, if, else, for, import, is, payable, override, virtual},
  keywordstyle=\color{codekw}\bfseries,
  comment=[l]{//}, morecomment=[s]{/*}{*/},
  commentstyle=\color{codecom}\itshape,
  stringstyle=\color{codestr},
  morestring=[b]"
}
\lstset{
  basicstyle=\ttfamily\small,
  backgroundcolor=\color{codebg},
  frame=single, rulecolor=\color{navy!20},
  numbers=left, numberstyle=\tiny\color{gray},
  breaklines=true, breakatwhitespace=false,
  columns=fullflexible, keepspaces=true,
  showstringspaces=false, tabsize=2,
  captionpos=b, xleftmargin=1.5em, framexleftmargin=1.5em,
  literate={à}{{\`a}}1 {è}{{\`e}}1 {é}{{\'e}}1 {ì}{{\`i}}1 {ò}{{\`o}}1 {ù}{{\`u}}1
}
\renewcommand{\lstlistingname}{Listato}
\renewcommand{\lstlistlistingname}{Elenco dei listati}

% ---------- Bibliografia ----------
\usepackage[backend=bibtex, style=numeric, sorting=none]{biblatex}
\addbibresource{bibliografia.bib}

% ---------- Collegamenti e URL ----------
\usepackage{url}
\usepackage[hidelinks]{hyperref}

% ---------- Metadati ----------
% ============================================================
%  METADATI DEL PROGETTO
%  Modifica SOLO questo file per aggiornare frontespizio,
%  intestazioni e informazioni del PDF.
% ============================================================

% --- Progetto ---
\newcommand{\titoloprogetto}{Certificati di sicurezza e SBOM\\per gli sviluppatori di software}
\newcommand{\titolobreve}{Certificati di sicurezza e SBOM} % usato nell'intestazione di pagina
\newcommand{\sottotitolo}{Progettazione di software sicuro e blockchain}

% --- Corso e ateneo ---
\newcommand{\corso}{Software Security e Blockchain}
\newcommand{\corsodilaurea}{Laurea Magistrale in Ingegneria Informatica e dell'Automazione}
\newcommand{\ateneo}{Università Politecnica delle Marche}
\newcommand{\dipartimento}{Dipartimento di Ingegneria dell'Informazione}
\newcommand{\annoaccademico}{2026/2027}

% --- Docenza (sostituisci con i nomi reali) ---
\newcommand{\docente}{Prof. Luca Spalazzi}
\newcommand{\assistente}{Dott. Tariq}

% --- Gruppo di lavoro (4 studenti): nome e matricola ---
\newcommand{\studenteA}{Nicola De Cinque}   \newcommand{\matricolaA}{S000001}
\newcommand{\studenteB}{Serena Stefano}   \newcommand{\matricolaB}{S000002}
\newcommand{\studenteC}{Alessandro Stella}   \newcommand{\matricolaC}{1137420}
\newcommand{\studenteD}{Gabriele Vianello}   \newcommand{\matricolaD}{S000004}

% --- Data di consegna ---
\newcommand{\dataconsegna}{\today}

% --- Repository del progetto (facoltativo, lascia vuoto per nasconderlo) ---
\newcommand{\repository}{https://github.com/BlockchainzGurus/eth-adam}

\hypersetup{
  pdftitle={\titolobreve},
  pdfauthor={\studenteA, \studenteB, \studenteC, \studenteD},
  pdfsubject={\corso - \ateneo},
  pdfkeywords={software security, blockchain, SBOM, certificati di sicurezza}
}

% ---------- Stile titoli ----------
\titleformat{\chapter}[hang]{\sffamily\bfseries\Huge\color{navy}}{\thechapter}{0.6em}{}[{\color{teal}\titlerule[1.5pt]}]
\titlespacing*{\chapter}{0pt}{-10pt}{20pt}
\titleformat{\section}{\sffamily\bfseries\Large\color{navy}}{\thesection}{0.6em}{}
\titleformat{\subsection}{\sffamily\bfseries\large\color{navy!85}}{\thesubsection}{0.6em}{}

% ---------- Intestazioni e piè di pagina ----------
\setlength{\headheight}{16pt}
\pagestyle{fancy}
\fancyhf{}
\fancyhead[L]{\sffamily\small\color{gray}\titolobreve}
\fancyhead[R]{\sffamily\small\color{gray}\leftmark}
\fancyfoot[C]{\sffamily\small\thepage}
\renewcommand{\headrulewidth}{0.4pt}
\renewcommand{\headrule}{\hbox to\headwidth{\color{teal}\leaders\hrule height \headrulewidth\hfill}}
\fancypagestyle{plain}{\fancyhf{}\fancyfoot[C]{\sffamily\small\thepage}\renewcommand{\headrulewidth}{0pt}}

% ---------- Comandi utili per il progetto ----------
\newcommand{\sbom}{\textsc{sbom}}
\newcommand{\todo}[1]{\textcolor{red}{\textbf{[TODO: #1]}}}
\newcommand{\codice}[1]{\texttt{#1}}
\newenvironment{riquadro}[1]{%
  \par\medskip\noindent\begin{tikzpicture}
  \node[draw=teal, fill=teal!6, rounded corners=3pt, line width=0.8pt,
        inner sep=8pt, text width=\dimexpr\linewidth-16pt-1.6pt\relax]\bgroup
  \sffamily\bfseries\color{navy}#1\par\smallskip\rmfamily\normalfont\color{black}%
}{\egroup;\end{tikzpicture}\par\medskip}

% ============================================================
\begin{document}

\pagenumbering{gobble}
% ============================================================
%  FRONTESPIZIO
%  Grafica: fascia navy con catena di blocchi (blockchain) e
%  scudo con spunta (sicurezza / certificato), codice a
%  barre di hash decorativo in fondo alla pagina.
%  I dati arrivano da metadati.tex: non serve modificarlo.
% ============================================================
\begin{titlepage}
\thispagestyle{empty}
\pagecolor{white}

% ---------- Grafica di sfondo ----------
\begin{tikzpicture}[remember picture, overlay]

  % Fascia superiore
  \fill[navy] (current page.north west) rectangle
              ([yshift=-6.2cm]current page.north east);
  \fill[teal] ([yshift=-6.2cm]current page.north west) rectangle
              ([yshift=-6.35cm]current page.north east);

  % Reticolo di punti (effetto "rete")
  \begin{scope}
    \clip (current page.north west) rectangle
          ([yshift=-6.2cm]current page.north east);
    \foreach \x in {0,0.8,...,21} {
      \foreach \y in {0,0.8,...,6.2} {
        \fill[white, opacity=0.06]
          ([xshift=\x cm, yshift=-\y cm]current page.north west) circle (0.03);
      }
    }
  \end{scope}

  % Intestazione di ateneo (testo nella fascia)
  \node[anchor=north west, text=white, align=left, inner sep=0pt]
    at ([xshift=2.2cm, yshift=-1.3cm]current page.north west)
    {\sffamily\bfseries\large \ateneo\\[2pt]
     \sffamily\small\color{white!80!navy} \dipartimento};

  % Logo opzionale (immagini/logo-univpm.png, sfondo bianco arrotondato)
  \IfFileExists{immagini/logo-univpm.png}{%
    \node[anchor=north east, fill=white, rounded corners=3pt, inner sep=4pt]
      at ([xshift=-2.2cm, yshift=-1.0cm]current page.north east)
      {\includegraphics[height=1.4cm]{immagini/logo-univpm.png}};
  }{%
    % Scudo con spunta se il logo non c'è
    \begin{scope}[shift={([xshift=-3.2cm, yshift=-1.9cm]current page.north east)}, scale=0.8]
      \path[fill=teal!85, draw=white, line width=1.2pt]
        (0,1) -- (0.95,0.7) -- (0.95,-0.1)
        .. controls (0.95,-0.75) and (0.3,-1.15) .. (0,-1.3)
        .. controls (-0.3,-1.15) and (-0.95,-0.75) .. (-0.95,-0.1)
        -- (-0.95,0.7) -- cycle;
      \draw[white, line width=2.2pt, line cap=round, line join=round]
        (-0.42,-0.02) -- (-0.1,-0.38) -- (0.45,0.38);
    \end{scope}
  }

  % Catena di blocchi (blockchain)
  \coordinate (o) at ([xshift=2.2cm, yshift=-4.1cm]current page.north west);
  \foreach \i/\h in {0/a3f9, 1/7c2e, 2/e41b, 3/09d8, 4/5bb6, 5/c7a1} {
    \draw[teal!80!white, line width=1pt, rounded corners=2pt, fill=white, fill opacity=0.07]
      ([xshift=2.75cm*\i]o) rectangle ++(1.5,-1.0);
    \node[font=\ttfamily\scriptsize, text=white, anchor=center]
      at ([xshift=2.75cm*\i+0.75cm, yshift=-0.32cm]o) {\#\h};
    \node[font=\ttfamily\tiny, text=teal!40!white, anchor=center]
      at ([xshift=2.75cm*\i+0.75cm, yshift=-0.7cm]o) {block \i};
  }
  \foreach \i in {0,...,4} {
    \draw[teal!80!white, line width=1pt, -{Stealth[length=4pt]}]
      ([xshift=2.75cm*\i+1.5cm, yshift=-0.5cm]o) -- ++(1.25,0);
  }

  % Striscia "codice a barre" di hash in fondo alla pagina
  \begin{scope}
    \pgfmathsetseed{42}
    \foreach \i in {0,...,59} {
      \pgfmathsetmacro{\w}{0.04 + 0.1*rnd}
      \pgfmathsetmacro{\o}{0.25 + 0.55*rnd}
      \fill[navy, opacity=\o]
        ([xshift=2.2cm+0.26cm*\i, yshift=1.3cm]current page.south west)
        rectangle ++(\w, 0.55);
    }
  \end{scope}
  \node[font=\ttfamily\tiny, text=navy!70, anchor=north west]
    at ([xshift=2.2cm, yshift=1.25cm]current page.south west)
    {sha256: 9f86d081884c7d659a2feaa0c55ad015a3bf4f1b2b0b822cd15d6c15b0f00a08};
\end{tikzpicture}

% ---------- Contenuto testuale ----------
\vspace*{4.6cm}

\begin{flushleft}
  {\sffamily\bfseries\color{teal}\large\MakeUppercase{\corso}}\\[3pt]
  {\sffamily\small\color{gray} \corsodilaurea}\\[1.4cm]

  {\sffamily\bfseries\fontsize{26}{30}\selectfont\color{navy} \titoloprogetto}\\[0.5cm]
  {\sffamily\large\color{navy!70} \sottotitolo}\\[0.3cm]
  {\color{teal}\rule{4cm}{2pt}}
\end{flushleft}

\vspace{1.2cm}

% Gruppo di lavoro
\begin{minipage}[t]{0.56\textwidth}
  {\sffamily\bfseries\small\color{teal}\MakeUppercase{Gruppo di lavoro}}\\[6pt]
  \begin{tabular}{@{}l@{\hspace{1.2em}}l@{}}
    \studenteA & {\ttfamily\small\color{gray}\matricolaA}\\[2pt]
    \studenteB & {\ttfamily\small\color{gray}\matricolaB}\\[2pt]
    \studenteC & {\ttfamily\small\color{gray}\matricolaC}\\[2pt]
    \studenteD & {\ttfamily\small\color{gray}\matricolaD}
  \end{tabular}
\end{minipage}%
\begin{minipage}[t]{0.40\textwidth}
  {\sffamily\bfseries\small\color{teal}\MakeUppercase{Docente}}\\[4pt]
  \docente\\[0.5cm]
  {\sffamily\bfseries\small\color{teal}\MakeUppercase{Assistente del docente}}\\[4pt]
  \assistente
\end{minipage}

\vfill

\begin{flushleft}
  {\sffamily\small Anno accademico \annoaccademico \quad\textbullet\quad \dataconsegna}
  \ifx\repository\empty\else\\[2pt]{\ttfamily\footnotesize\color{gray}\url{\repository}}\fi
\end{flushleft}
\vspace{2.4cm}
\end{titlepage}


\pagenumbering{roman}

% ---------- Abstract ----------
\chapter*{Abstract}
\addcontentsline{toc}{chapter}{Abstract}
% Sintesi di 150-250 parole: problema, approccio, risultati principali.
Questo progetto affronta la progettazione di software sicuro lungo tutta la catena di fornitura,
combinando \emph{certificati di sicurezza} e \emph{Software Bill of Materials} (\sbom) per fornire
agli sviluppatori garanzie verificabili sull'origine e sull'integrità dei componenti utilizzati.
La tecnologia blockchain è impiegata come registro non modificabile per l'ancoraggio dei
certificati e degli hash delle \sbom.

\todo{Completare l'abstract con obiettivi, metodo e risultati ottenuti.}


\tableofcontents
\listoffigures
\listoftables
\lstlistoflistings

% ---------- Corpo del documento ----------
\clearpage
\pagenumbering{arabic}

\chapter{Introduzione}
\label{cap:introduzione}

\section{Contesto e motivazione}
\todo{Perché la sicurezza della supply chain software è un problema attuale; attacchi noti, normative e linee guida di riferimento.}

\section{Obiettivi del progetto}
\todo{Elencare gli obiettivi: certificati di sicurezza per gli sviluppatori, generazione e verifica della \sbom, ruolo della blockchain.}

\section{Contributi}
\todo{Descrivere in modo sintetico i risultati originali del gruppo.}

\section{Struttura della relazione}
Il Capitolo~\ref{cap:stato-arte} presenta lo stato dell'arte, il Capitolo~\ref{cap:requisiti}
definisce requisiti e modello delle minacce, il Capitolo~\ref{cap:progettazione} descrive la progettazione,
il Capitolo~\ref{cap:implementazione} l'implementazione, il Capitolo~\ref{cap:valutazione} la valutazione
e il Capitolo~\ref{cap:conclusioni} riporta le conclusioni.

\chapter{Stato dell'arte}
\label{cap:stato-arte}

\section{Sicurezza della supply chain del software}
\todo{Pratiche di sviluppo sicuro e riferimenti normativi, ad esempio \cite{nist-ssdf}.}

\section{Software Bill of Materials}
Una \sbom{} è un inventario strutturato dei componenti che compongono un software, con relazioni di
dipendenza, versioni e licenze. I formati più diffusi sono \emph{SPDX}~\cite{spdx} e \emph{CycloneDX}~\cite{cyclonedx}.

\begin{table}[htbp]
  \centering
  \caption{Confronto tra formati \sbom{} (da completare).}
  \label{tab:formati-sbom}
  \begin{tabularx}{\linewidth}{@{}lXX@{}}
    \toprule
    \textbf{Aspetto} & \textbf{SPDX} & \textbf{CycloneDX}\\
    \midrule
    Ambito principale & \todo{} & \todo{}\\
    Formati di serializzazione & \todo{} & \todo{}\\
    Supporto a vulnerabilità & \todo{} & \todo{}\\
    \bottomrule
  \end{tabularx}
\end{table}

\section{Certificati di sicurezza e firma del software}
\todo{Certificati X.509, firma del codice, attestazioni di provenienza e infrastrutture a chiave pubblica.}

\section{Blockchain per l'integrità e la tracciabilità}
\todo{Registro distribuito, smart contract, ancoraggio di hash~\cite{nakamoto}.}

\chapter{Requisiti e modello delle minacce}
\label{cap:requisiti}

\section{Attori e casi d'uso}
\todo{Sviluppatore, verificatore, ente certificatore, consumatore del software.}

\section{Requisiti funzionali e non funzionali}
\begin{table}[htbp]
  \centering
  \caption{Elenco dei requisiti.}
  \label{tab:requisiti}
  \begin{tabularx}{\linewidth}{@{}llX@{}}
    \toprule
    \textbf{ID} & \textbf{Tipo} & \textbf{Descrizione}\\
    \midrule
    RF-01 & Funzionale & \todo{Generazione della \sbom{} a partire dal progetto.}\\
    RF-02 & Funzionale & \todo{Emissione e verifica del certificato di sicurezza.}\\
    RNF-01 & Sicurezza & \todo{Integrità e non ripudio dei certificati.}\\
    \bottomrule
  \end{tabularx}
\end{table}

\section{Modello delle minacce}
\todo{Metodologia adottata (ad esempio STRIDE), asset, avversari, superfici di attacco e contromisure.}

\chapter{Progettazione}
\label{cap:progettazione}

\section{Architettura generale}
La Figura~\ref{fig:architettura} mostra uno schema d'esempio dell'architettura; va adattato al progetto.

\begin{figure}[htbp]
  \centering
  \begin{tikzpicture}[
      node distance=1.4cm and 1.6cm,
      box/.style={draw=navy, fill=navy!6, rounded corners=3pt, minimum width=2.8cm,
                  minimum height=1cm, align=center, font=\sffamily\small},
      chain/.style={draw=teal, fill=teal!10, rounded corners=3pt, minimum width=2.8cm,
                    minimum height=1cm, align=center, font=\sffamily\small},
      arr/.style={-{Stealth[length=4pt]}, thick, color=navy}]
    \node[box] (dev) {Sviluppatore};
    \node[box, right=of dev] (sbom) {Generazione\\\sbom};
    \node[box, right=of sbom] (cert) {Emissione\\certificato};
    \node[chain, below=of cert] (bc) {Blockchain\\(smart contract)};
    \node[box, below=of dev] (ver) {Verificatore};
    \draw[arr] (dev) -- (sbom);
    \draw[arr] (sbom) -- (cert);
    \draw[arr] (cert) -- node[right, font=\scriptsize]{hash} (bc);
    \draw[arr] (ver) -- node[above, font=\scriptsize]{verifica} (bc);
  \end{tikzpicture}
  \caption{Architettura di riferimento (esempio da adattare).}
  \label{fig:architettura}
\end{figure}

\section{Struttura della \sbom}
\todo{Campi, formato scelto e motivazioni.}

\section{Certificati di sicurezza}
\todo{Contenuto del certificato, catena di fiducia, ciclo di vita (emissione, revoca, scadenza).}

\section{Ruolo della blockchain}
\todo{Dati ancorati on-chain, scelta della rete, costi e compromessi.}

\chapter{Implementazione}
\label{cap:implementazione}

\section{Ambiente e strumenti}
\todo{Linguaggi, librerie, framework di sviluppo e test, rete di test.}

\section{Smart contract}
Il Listato~\ref{lst:registro} riporta un esempio minimale di registro di hash, da sostituire con il codice reale.

\begin{lstlisting}[language=Solidity, caption={Esempio: registro di hash di certificati e \sbom.}, label={lst:registro}]
// SPDX-License-Identifier: MIT
pragma solidity ^0.8.20;

contract RegistroCertificati {
    mapping(bytes32 => address) public emittente;

    event Registrato(bytes32 indexed hash, address indexed da);

    function registra(bytes32 hash) public {
        require(emittente[hash] == address(0), "gia' registrato");
        emittente[hash] = msg.sender;
        emit Registrato(hash, msg.sender);
    }
}
\end{lstlisting}

\section{Generazione della \sbom}
\begin{lstlisting}[language=bash, caption={Esempio di comando per generare una SBOM.}, label={lst:cmd-sbom}]
# esempio generico: sostituire con lo strumento realmente utilizzato
cyclonedx-py environment -o sbom.json
\end{lstlisting}

\section{Emissione e verifica dei certificati}
\todo{Descrivere i flussi implementati e le scelte crittografiche.}

\chapter{Valutazione}
\label{cap:valutazione}

\section{Metodologia di test}
\todo{Casi di test, strumenti di analisi statica e dinamica, ambiente.}

\section{Analisi di sicurezza}
\todo{Verifica delle contromisure rispetto al modello delle minacce (Capitolo~\ref{cap:requisiti}).}

\section{Costi e prestazioni}
\todo{Costo del gas, tempi di generazione e verifica, scalabilità.}

\section{Limiti della soluzione}
\todo{Limiti noti e lavori futuri.}

\chapter{Conclusioni}
\label{cap:conclusioni}

\todo{Riepilogo dei risultati rispetto agli obiettivi dell'Introduzione e possibili sviluppi futuri.}


% ---------- Appendici ----------
\appendix
\chapter{Ripartizione del lavoro}
\label{app:ripartizione}

\begin{table}[htbp]
  \centering
  \caption{Contributo dei componenti del gruppo.}
  \label{tab:ripartizione}
  \begin{tabularx}{\linewidth}{@{}lX@{}}
    \toprule
    \textbf{Studente} & \textbf{Attività svolte}\\
    \midrule
    \studenteA & \todo{}\\
    \studenteB & \todo{}\\
    \studenteC & \todo{}\\
    \studenteD & \todo{}\\
    \bottomrule
  \end{tabularx}
\end{table}


% ---------- Bibliografia ----------
\clearpage
\addcontentsline{toc}{chapter}{Bibliografia}
\printbibliography[heading=none]

\end{document}
